\documentclass[11pt,a4paper]{article}
\usepackage[T1]{fontenc}
\usepackage{lmodern}
\usepackage[margin=25mm,headheight=15pt]{geometry}
\usepackage{amsmath,amssymb,amsthm,mathtools,mathrsfs}
\usepackage{microtype}
\usepackage{booktabs,array,longtable}
\usepackage{xcolor}
\usepackage{enumitem}
\usepackage{fancyhdr}
\usepackage{needspace}
\usepackage{listings}
\usepackage{hyperref}
\usepackage{bookmark}
\definecolor{ink}{HTML}{173D4F}
\hypersetup{colorlinks=true,linkcolor=ink,citecolor=ink,urlcolor=ink,
 pdftitle={Polynomial Periods and Sharp Residue Hierarchies for the Global Inverse Fueter Map},
 pdfauthor={Research article prepared for Vladimir Reshetnikov},
 pdfsubject={Hermite interpolation, global Fueter inversion, polynomial monodromy and sharp singularity bounds}}
\pagestyle{fancy}
\fancyhf{}
\fancyhead[L]{\small Polynomial periods for the inverse Fueter map}
\fancyhead[R]{\small Research article}
\fancyfoot[C]{\thepage}
\setlength{\parindent}{1.2em}
\setlength{\parskip}{3pt}
\setlength{\emergencystretch}{2em}
\numberwithin{equation}{section}
\newtheorem{theorem}{Theorem}[section]
\newtheorem{lemma}[theorem]{Lemma}
\newtheorem{proposition}[theorem]{Proposition}
\newtheorem{corollary}[theorem]{Corollary}
\theoremstyle{definition}
\newtheorem{definition}[theorem]{Definition}
\theoremstyle{remark}
\newtheorem{remark}[theorem]{Remark}
\newcommand{\R}{\mathbb R}
\newcommand{\C}{\mathbb C}
\newcommand{\Z}{\mathbb Z}
\newcommand{\HH}{\mathbb H}
\newcommand{\ii}{\mathrm i}
\newcommand{\Cl}{\operatorname{Cl}}
\newcommand{\Hol}{\mathcal O}
\newcommand{\AM}{\mathcal A_h}
\newcommand{\Fu}{\mathcal F_h}
\newcommand{\PP}{\mathscr P_{<2h}(E)}
\newcommand{\PPl}[1]{\mathscr P_{<#1}(E)}
\newcommand{\Om}{\Omega_h}
\newcommand{\Per}{\operatorname{Per}}
\newcommand{\ord}{\operatorname{ord}}
\newcommand{\Span}{\operatorname{span}}
\newcommand{\dd}{\,\mathrm d}
\newcommand{\norm}[1]{\left\lVert#1\right\rVert}
\newcommand{\barz}{\overline z}
\newcommand{\barp}{\overline p}
\newcommand{\Rop}{\mathsf R}
\newcommand{\Sop}{\mathsf S}
\newcommand{\sig}{\sigma_{2h}}
\newcommand{\disc}{\mathbb D}
\newcommand{\res}{\operatorname{res}}
\newcommand{\sourcepath}{\texttt{SetTheory/Cardinals/docs/reports/}}
\lstset{basicstyle=\ttfamily\small,breaklines=true,columns=fullflexible,
 showstringspaces=false,frame=single,rulecolor=\color{black!25},
 xleftmargin=5pt,xrightmargin=5pt}

\begin{document}
\hypersetup{pageanchor=false}
\begin{titlepage}
\vspace*{12mm}
\noindent{\color{ink}\large ANALYSIS / GLOBAL INVERSION / CLIFFORD MODULES}
\vspace{13mm}
\begin{flushleft}
{\color{ink}\LARGE\bfseries Polynomial Periods and\\[5pt]
Sharp Residue Hierarchies\\[5pt]
for the Global Inverse Fueter Map}\par
\vspace{8mm}
{\large A Hermite-interpolation construction in every\\
even ambient dimension at least four}\par
\vspace{13mm}
Research article prepared for Vladimir Reshetnikov\\
September 2026
\end{flushleft}
\vspace{8mm}
\begin{abstract}
The quaternionic inverse Fueter problem has a two-period obstruction: a
single-valued slice-regular primitive exists precisely when two target-derived
closed one-forms have zero periods. We construct its higher-dimensional
counterpart. In ambient dimension $2h+2$, one explicit polynomial-valued
closed form of degree $2h-1$ captures all obstructions. Its normalization and
completeness follow from a two-node Hermite-interpolation identity, proved for
arbitrary holomorphic stems by an exact finite-jet argument. Integrating its
coefficients and evaluating the resulting polynomial at the moving point
constructs every global inverse.

The period polynomial gives complete monodromy, an explicit finite-connectivity
cokernel, and the smallest covering on which inversion is possible. Near a
punctured conjugacy sphere it has a finer structure: divisibility by a real
quadratic determines the least possible singularity order among all targets
with that residue. We prove sharp axial-pair growth bounds and sharp physical
$L^2$ energy constants, exhibit logarithmic representatives attaining them,
and classify the residue classes compatible with local $L^p$ integrability
for $p\ge1$. The arguments are self-contained within smooth axial analysis on
product domains disjoint from the real axis.
\end{abstract}
\vfill
\noindent\small\textbf{Status.} These are written mathematical proofs and proposed
extensions of the inspected ProveIt project, not a claim of established priority
or a resolution of an identified published conjecture. The article is unrefereed
and is not formally verified in Lean. The accompanying finite symbolic and
numerical checks support the derivations but do not replace them.
\end{titlepage}
\hypersetup{pageanchor=true}
\pagenumbering{roman}

\setcounter{tocdepth}{2}
\begingroup
\small
\setlength{\parskip}{0pt}
\tableofcontents
\endgroup
\clearpage
\pagenumbering{arabic}

\section{The questions and the contribution}
\label{sec:intro}

Local Fueter inversion, global single-valued inversion, and the singularity
forced by a global obstruction are distinct questions. A local inverse can
change after continuation around a hole; a quantitative theory must also
measure the cost of that change.

The inspected ProveIt report on slice regularity and Fueter inversion
\cite{repo} gives a quaternionic pair of target-derived closed forms, affine
monodromy, an $8b$-dimensional real cokernel for $b$ holes, and a distinction
between axial-pair and whole-sphere norms. We rederive that starting point and
extend it to higher dimensions, including a residue filtration with optimal
constants.

The questions are concrete: is there a finite target-only period system in
every dimension covered by the differential Fueter map? Are its periods
complete and freely realizable? Do they determine a sharp minimum singularity
scale? The answers below are affirmative under explicit hypotheses.

\subsection{Main results in one place}

Fix $h\ge1$, a connected planar domain $D\subset\{\operatorname{Im}z>0\}$,
and a real coefficient module $E$ as in Section~\ref{sec:setting}. The physical
ambient dimension is $2h+2$, and $d=\dim_{\R}E$. Write an axially monogenic target
as $g=P+IQ$, and put
\begin{equation}\label{eq:intro-current}
\begin{split}
\Om[g](t)={}&\kappa_h\bigl((t-x)^2+y^2\bigr)^{h-1}\\
&\quad\cdot\left[\bigl((x-t)P+yQ\bigr)\dd x
 +\bigl(yP-(x-t)Q\bigr)\dd y\right],\\
\kappa_h={}&\frac{(-1)^{h+1}}{2^{2h-1}h!(h-1)!}.
\end{split}
\end{equation}
Here $t$ is a real formal variable, not a coordinate on $D$.

\begin{theorem}[Summary of the global and local conclusions]\label{thm:summary}
With the definitions and hypotheses made precise in the subsequent sections:
\begin{enumerate}[label=\textup{(\roman*)},leftmargin=2em]
\item The form \eqref{eq:intro-current} is closed coefficientwise. A global
holomorphic stem $F$ satisfying $\Fu F=g$ exists if and only if all its period
polynomials vanish. Every inverse is obtained by coefficientwise integration
followed by evaluation at $t=z$.
\item If $D$ has $b$ holes with the regularity specified in
Section~\ref{sec:cokernel}, then the real cokernel is isomorphic to
$\PP^b$ and has dimension $2hdb$. Every period assignment has an explicit
logarithmic representative.
\item Let $M\ne0$ be the residue polynomial at $p=u+\ii v$, $v>0$, and let
$r=\ord_p M$. Then $0\le r\le h-1$, and the least possible pole-growth
order of a target with residue $M$ is $s=h-r$. Equivalently, the residues
realizable with growth $O(|z-p|^{-s})$ are exactly
\[
 \bigl((t-u)^2+v^2\bigr)^{h-s}\PPl{2s},\qquad 0\le s\le h.
\]
The level $s$ has real dimension $2sd$.
\item The least leading axial-pair size is
\[
 A_{h,p,M}=\frac{2^h h!(s-1)!}{2\pi v^h}\norm{M^{(r)}(p)}.
\]
Every representative satisfies the corresponding lower bound, and an explicit
one attains it. Its physical energy is logarithmically divergent at $s=1$ and
has optimal power divergence $\varepsilon^{2-2s}$ for $s\ge2$; the exact
constants are given in Theorem~\ref{thm:energy}.
\end{enumerate}
\end{theorem}

The displayed conclusions are proved below, not inferred from a finite
experiment. The summary's reference to a \emph{least} growth order is important:
a target may have extra zero-residue singular terms and therefore grow faster
than the obstruction itself requires.

\subsection{What is inherited, what is extended, and what is not claimed}

The Fueter--Sce mapping, local inverse formulas, and a finite-dimensional
polynomial kernel belong to established theory
\cite{colombo2014,colombo2013,demartino2023}. We reprove the particular radial
formula and kernel needed here to fix all constants and hypotheses. Their
existence is not presented as a new result.

The proposed additions are the normalized target-only polynomial current,
the moving Hermite-polynomial identity that makes it complete, the constructive
period splitting, and the higher-dimensional residue filtration with sharp
size and energy constants. These provide explicit answers to the extension
questions posed above. A focused comparison with the cited primary sources
and the inspected repository report did not establish priority for these
formulas. Absence from that comparison is not proof of absence from the wider
literature.

There is also no claimed contradiction with published surjectivity results.
For example, Perotti's Theorem~2 explicitly assumes that every connected
component of the planar stem domain is simply connected \cite{perotti2024}.
Our quaternionic criterion is compatible with that theorem; on upper-half-plane
product domains it explains precisely what can obstruct the same conclusion
when that topological hypothesis is removed.

\section{Setting, coefficients, and the two complex structures}
\label{sec:setting}

\subsection{An orthogonal Clifford module}

Let $m=2h+1$. Let $E$ be a finite-dimensional real Hilbert space carrying
operators $J_1,\ldots,J_m$ such that
\begin{equation}\label{eq:clifford}
 J_j^*=-J_j,\qquad J_jJ_k+J_kJ_j=-2\delta_{jk}\operatorname{Id}_E.
\end{equation}
Thus $E$ is an orthogonal module for $\Cl_{0,m}$. For
$I=(I_1,\ldots,I_m)\in S^{m-1}$, write $J_I=\sum I_jJ_j$ and abbreviate
$J_I a$ to $Ia$. Then $J_I^2=-\operatorname{Id}_E$ and $J_I$ is an isometry.

The regular Clifford module is one example. For $h=1$, the quaternionic
module $E=\HH$, with the three operators given by left multiplication by the
quaternionic imaginary units, is another. In the latter case $d=4$; this is
not an identification of the eight-dimensional algebra $\Cl_{0,3}$ with $\HH$.

Let $E_{\C}=E\oplus\ii E$ be the complexification of $E$. The symbol $\ii$ is
\emph{central}; it is not any one of the operators $J_I$. Complex conjugation
acts by $\overline{a+\ii b}=a-\ii b$, and
\[
 \norm{a+\ii b}^2=\norm a^2+\norm b^2.
\]
All polynomial coefficients below lie in the real space $E$, unless the
contrary is explicitly stated. This distinction is essential to the kernel.

\subsection{Stems and axial targets}

Let $D$ be a connected open subset of $\C^+=\{x+\ii y:y>0\}$. Its circularization
is
\[
 U_D=\{(x,yI)\in\R\times\R^m:x+\ii y\in D,\ I\in S^{m-1}\}.
\]
We use $q=x+yI$ as coordinate notation for this point; no multiplication of
points in $\R^{m+1}$ is needed. A holomorphic map
\[
 F=A+\ii B\in\Hol(D;E_{\C})
\]
induces the slice function $f(q)=A(x,y)+IB(x,y)$. Holomorphicity means
\begin{equation}\label{eq:CR}
 A_x=B_y,\qquad A_y=-B_x.
\end{equation}
This upper-half-plane formulation is equivalent to defining a conjugate-symmetric
stem on the disjoint union $D\cup\overline D$. It imposes no matching condition
across the real axis because that axis is excluded.

\begin{definition}[Axially monogenic target]\label{def:AM}
The real vector space $\AM(D;E)$ consists of smooth pairs $(P,Q):D\to E^2$
that satisfy the Vekua system
\begin{equation}\label{eq:vekua}
 P_x-Q_y=\frac{2h}{y}Q,\qquad P_y+Q_x=0.
\end{equation}
We identify the pair with $g(x+yI)=P(x,y)+IQ(x,y)$ and set
$W=P+\ii Q\in E_{\C}$.
\end{definition}

Indeed, applying the Dirac operator
$\partial_x+\sum_{j=1}^mJ_j\partial_{x_j}$ to the axial expression gives
\[
 \left(P_x-Q_y-\frac{m-1}{y}Q\right)+I(P_y+Q_x).
\]
Vanishing for both $I$ and $-I$ is equivalent to \eqref{eq:vekua}. This also
shows that the coefficient formalism describes genuine monogenic functions
on $U_D$.

The norm $\norm W=(\norm P^2+\norm Q^2)^{1/2}$ is called the
\emph{axial-pair norm}. It is not in general equal to
$\max_{I\in S^{m-1}}\norm{P+IQ}$. The distinction will be retained throughout
our sharp estimates.

\subsection{The map and its normalization}

The Fueter map in this article is the unnormalized differential operator
\begin{equation}\label{eq:fueter-definition}
 \Fu F=\Delta_{2h+2}^{\,h}f,
 \qquad f=\mathcal I(F)=A+IB.
\end{equation}
It is real-linear. Multiplying the stem by the central scalar $\ii$ does not,
in general, multiply its image by an analogous scalar. Put
\begin{equation}\label{eq:constants}
 c_h=2^hh!,\qquad C_h=(-1)^h c_h^2,\qquad
 \kappa_h=\frac{(-1)^{h+1}}{2^{2h-1}h!(h-1)!}.
\end{equation}
The normalization identity used repeatedly below is
\begin{equation}\label{eq:calibration-constant}
 \kappa_h C_h=-2h.
\end{equation}
For example, $(c_1,\kappa_1)=(2,1/2)$ and
$(c_2,\kappa_2)=(8,-1/16)$.

\section{Radial calculus and two calibration monomials}
\label{sec:radial}

Define two operators on smooth functions of $(x,y)$, with $y>0$, by
\[
 \Rop u=\frac1y\partial_yu,\qquad
 \Sop u=\partial_y\!\left(\frac uy\right).
\]
The operators act componentwise on $E$-valued functions.

\begin{lemma}[Radial expression for the Fueter map]\label{lem:radial}
For every holomorphic stem $F=A+\ii B$,
\begin{equation}\label{eq:radial}
 \Fu F=P+IQ,\qquad
 P=c_h\Rop^hA,\qquad Q=c_h\Sop^hB.
\end{equation}
The resulting pair satisfies \eqref{eq:vekua}.
\end{lemma}

\begin{proof}
On a scalar axial component the ambient Laplacian acts as $L_{2h}$, where
\[
 L_s=\partial_x^2+\partial_y^2+\frac{s}{y}\partial_y.
\]
On $IB$ it acts as $I M_{2h}B$, where
\[
 M_s=\partial_x^2+\partial_y^2+\frac{s}{y}\partial_y-\frac{s}{y^2}.
\]
For arbitrary real $s$, direct differentiation gives
\[
 L_s\Rop=\Rop L_{s-2},\qquad M_s\Sop=\Sop M_{s-2}.
\]
Because $A$ and $B$ are harmonic in the planar variables,
$L_sA=s\Rop A$ and $M_sB=s\Sop B$. An induction therefore gives
\[
 \Delta_{2h+2}^{\,j}(A+IB)
 =\prod_{a=0}^{j-1}(2h-2a)\,\bigl(\Rop^jA+I\Sop^jB\bigr)
 \quad(0\le j\le h).
\]
At $j=h$ the product is $c_h$.

The identities
\[
 \partial_y\Rop^h=\Sop^h\partial_y,
 \qquad
 \Rop^h\partial_y=\left(\partial_y+\frac{2h}{y}\right)\Sop^h
\]
follow by induction from the definitions. Consequently
\[
 P_y+Q_x=c_h\Sop^h(A_y+B_x)=0,
\]
and
\[
 P_x=c_h\Rop^h B_y
 =c_h\left(\partial_y+\frac{2h}{y}\right)\Sop^hB
 =Q_y+\frac{2h}{y}Q.
\]
These are the two required equations.
\end{proof}

\begin{lemma}[Calibration]\label{lem:calibration}
For a real coefficient $a\in E$ and integer $k\ge0$,
\begin{align}
 \Fu(z^ka)&=0 &&(k<2h),\label{eq:low-calibration}\\
 \Fu(z^{2h}a)&=C_h a,\label{eq:even-calibration}\\
 \Fu(z^{2h+1}a)&=C_h\bigl((2h+1)xa+Iy a\bigr).
 \label{eq:odd-calibration}
\end{align}
\end{lemma}

\begin{proof}
The slice lift of $z^ka$ is a polynomial of total degree $k$ in the ambient
coordinates. In fact its even powers of $y$ are powers of
$x_1^2+\cdots+x_m^2$, and its odd terms have $Iy=\sum x_jJ_j$ as a factor.
An operator of differential order $2h$ annihilates it when $k<2h$.

For $z^{2h}$ the only real-part term surviving $\Rop^h$ is
$(-1)^hy^{2h}$, and
$\Rop^hy^{2h}=2^hh!=c_h$. All imaginary terms are annihilated by $\Sop^h$.
For $z^{2h+1}$, the relevant highest terms are
$(-1)^h(2h+1)xy^{2h}$ in the real part and
$(-1)^hy^{2h+1}$ in the imaginary part. Since
$\Sop^hy^{2h+1}=c_hy$, the remaining formulas follow.
\end{proof}

\begin{remark}[Why the coefficient field matters]
For $h=1$, the constant stem $F=\ii a$ has $A=0$ and $B=a$, whence
$\mathcal F_1F=-2Ia/y^2$. It is not in the kernel unless $a=0$.
Thus a claim that the kernel consists of arbitrary \emph{complex}-coefficient
polynomials of degree below $2h$ would already be false for constants.
\end{remark}

\section{A target-only polynomial current}
\label{sec:current}

For a real formal variable $t$, let
\begin{equation}\label{eq:rt}
 R_t=(t-x)^2+y^2=(t-z)(t-\barz).
\end{equation}
Let $\PPl{k}$ denote the real vector space of $E$-coefficient polynomials in
$t$ of degree less than $k$, with $\PPl{0}=\{0\}$. Its real dimension is $kd$.
The current $\Om[g]$ defined in \eqref{eq:intro-current} takes values in $\PP$:
its degree in $t$ is at most $2h-1$.

\begin{theorem}[Closedness]\label{thm:closed}
For every $g\in\AM(D;E)$, $\mathrm d\Om[g]=0$ coefficientwise.
\end{theorem}

\begin{proof}
Set $a=x-t$, $L=aP+yQ$, and $N=yP-aQ$. The Vekua system yields
\begin{align*}
 N_x-L_y
 &=y(P_x-Q_y)-2Q-a(Q_x+P_y)\\
 &=2(h-1)Q.
\end{align*}
For $h>1$, the contribution from differentiating the factor $R_t^{h-1}$ is
\begin{align*}
 (\partial_xR_t^{h-1})N-(\partial_yR_t^{h-1})L
 &=2(h-1)R_t^{h-2}(aN-yL)\\
 &=-2(h-1)R_t^{h-1}Q.
\end{align*}
It cancels the preceding contribution. For $h=1$ the factor is constant and
both contributions are zero. Multiplication by $\kappa_h$ preserves closedness.
\end{proof}

A useful equivalent expression uses central complex notation:
\begin{equation}\label{eq:complex-current}
\begin{split}
 \Om[g](t)=-\frac{\kappa_h}{2}\bigl[
 &(t-z)^{h-1}(t-\barz)^h W\dd z\\
 {}+{}&(t-z)^h(t-\barz)^{h-1}\overline W\dd\barz\bigr].
\end{split}
\end{equation}
To verify it, write the real-form coefficients as $L,N$. Then
$L-\ii N=(\barz-t)W$ and
$L+\ii N=(z-t)\overline W$, and use
$\dd x=(\dd z+\dd\barz)/2$ and
$\dd y=(\dd z-\dd\barz)/(2\ii)$.

\begin{lemma}[Injectivity of the current]\label{lem:current-injective}
If $\Om[g_1]=\Om[g_2]$, then $g_1=g_2$.
\end{lemma}

\begin{proof}
Apply the assertion to the difference and fix $(x,y)$. Evaluate the polynomial
identity at the real number $t=x$. Its $\dd x$ coefficient is
$\kappa_h y^{2h-1}Q$ and its $\dd y$ coefficient is
$\kappa_h y^{2h-1}P$. Neither $\kappa_h$ nor $y$ is zero.
\end{proof}

\subsection{Recovery of the quaternionic pair}

When $h=1$,
\[
 \Om[g](t)=\theta_g+t\omega_g,
\]
where
\[
 \omega_g=\tfrac12(-P\dd x+Q\dd y),\qquad
 \theta_g=\tfrac12\bigl((xP+yQ)\dd x+(yP-xQ)\dd y\bigr).
\]
Thus the polynomial current specializes exactly to the two forms in the
inspected repository report \cite{repo}. It does not introduce a different
normalization of their periods.

\section{The moving Hermite-polynomial identity}
\label{sec:hermite}

The main completeness mechanism is not a repeated choice of harmonic
conjugates. It is a polynomial interpolation problem with two moving nodes.

\begin{definition}[Hermite polynomial of a stem]\label{def:hermite}
For $F\in\Hol(D;E_{\C})$ and $z\in D$, let $H_F(z;t)\in\PP$ be the unique
polynomial such that
\begin{align}
 \partial_t^j H_F(z;z)&=F^{(j)}(z),\label{eq:hermite-jets}\\
 \partial_t^j H_F(z;\barz)&=\overline{F^{(j)}(z)},
 \qquad 0\le j<h.\nonumber
\end{align}
\end{definition}

The nodes $z$ and $\barz$ are distinct. Componentwise polynomial Hermite
interpolation therefore gives a unique $E_{\C}$-coefficient polynomial of
degree below $2h$. Conjugating its coefficients and interchanging the nodes
preserves the interpolation data, so uniqueness makes its coefficients real,
i.e. elements of $E$. Its coefficients depend smoothly on $(x,y)$: the inverse
interpolation matrix is rational in the nodes with nonzero powers of
$z-\barz$ in its denominator.

\begin{theorem}[Hermite differential identity]\label{thm:hermite}
For every holomorphic stem $F$,
\begin{equation}\label{eq:hermite-identity}
 \mathrm d_{x,y}H_F(z;t)=\Om[\Fu F](t).
\end{equation}
The formal variable $t$ is held fixed while the left side is differentiated.
\end{theorem}

\begin{proof}
Fix a point $z_0\in D$. Both sides at $z_0$ are real-linear functions of the
finite jet
\[
 \bigl(F(z_0),F'(z_0),\ldots,F^{(h)}(z_0)\bigr).
\]
For the left side, differentiating the coefficients of the interpolation
problem introduces derivatives only through order $h$. For the right side,
this follows from the radial formula \eqref{eq:radial} and holomorphicity.

Every such jet is the jet at $z_0$ of a polynomial with coefficients in $E$
and degree at most $2h+1$. Indeed, interpolate the jet through order $h$
at $z_0$ and its conjugate jet at $\bar z_0$. The two-node Hermite problem
has $2h+2$ conditions, and its unique solution has real coefficients by the
same conjugation argument used above. Therefore it is enough to check
\eqref{eq:hermite-identity} for the $2h+2$ monomials
$1,z,\ldots,z^{2h+1}$ with an arbitrary coefficient in $E$.
This is an exact spanning argument for arbitrary jets, not a finite-range
experimental inference.

For $k<2h$, the interpolant of $z^k$ is the constant-in-$z$ polynomial $t^k$;
both sides vanish by Lemma~\ref{lem:calibration}.
For the remaining two monomials, put $R_t=(t-z)(t-\barz)$. Polynomial division
gives their Hermite remainders:
\begin{align}
 H_{z^{2h}}(z;t)&=t^{2h}-R_t^h,\label{eq:hermite-even}\\
 H_{z^{2h+1}}(z;t)&=t^{2h+1}-(t+2hx)R_t^h.
 \label{eq:hermite-odd}
\end{align}
The remainders have degree below $2h$ and differ from the original monomial
by a multiple of both $(t-z)^h$ and $(t-\barz)^h$, proving the asserted
interpolation conditions.

Differentiating \eqref{eq:hermite-even} gives
\[
 \mathrm dH_{z^{2h}}=-2hR_t^{h-1}\bigl((x-t)\dd x+y\dd y\bigr).
\]
This is the current of the constant target $C_h$, because
$\kappa_hC_h=-2h$. Differentiating \eqref{eq:hermite-odd} gives
\begin{equation*}
\begin{split}
 \mathrm dH_{z^{2h+1}}=-2hR_t^{h-1}\bigl[
 &\bigl((2h+1)x(x-t)+y^2\bigr)\dd x\\
 &+y(t+2hx)\dd y\bigr].
\end{split}
\end{equation*}
The target in \eqref{eq:odd-calibration} has
$P=C_h(2h+1)x$ and $Q=C_hy$, and substitution into its current gives exactly
this expression. These checks establish the identity for every jet at the
arbitrary point $z_0$, hence everywhere on $D$.
\end{proof}

\begin{corollary}[Polynomial kernel]\label{cor:kernel}
For connected $D\subset\C^+$,
\begin{equation}\label{eq:kernel}
 \ker\Fu=\{F(z)=K(z):K\in\PP\}.
\end{equation}
In particular, the kernel has real dimension $2hd$.
\end{corollary}

\begin{proof}
If $\Fu F=0$, Theorem~\ref{thm:hermite} says that every coefficient of
$H_F$ has zero differential. It is therefore constant on the connected domain:
$H_F(z;t)=K(t)$. Evaluating at $t=z$ gives $F(z)=K(z)$, with $K\in\PP$.
Conversely, Lemma~\ref{lem:calibration} annihilates every such polynomial.
\end{proof}

The existence of a polynomial kernel is classical
\cite{colombo2014,demartino2023}. The proof above is included because it identifies
precisely the coefficient space needed for the global obstruction and because
it uses the same identity as the inverse construction.

\section{A complete constructive global inverse theorem}
\label{sec:inverse}

For a piecewise smooth closed curve $\gamma$ in $D$, define its period
polynomial by
\begin{equation}\label{eq:period}
 \Per_g(\gamma;t)=\int_\gamma\Om[g](t)\in\PP.
\end{equation}
Integration is coefficientwise. Since the current is closed, its periods
are invariant under homotopy and depend only on the homology class of the
curve. The choice of a basepoint affects neither a closed-curve integral nor
its polynomial value.

\begin{theorem}[Global inverse and all solutions]\label{thm:inverse}
Let $g\in\AM(D;E)$. The following are equivalent:
\begin{enumerate}[label=\textup{(\alph*)},leftmargin=2em]
\item There exists $F\in\Hol(D;E_{\C})$ such that $\Fu F=g$.
\item $\Per_g(\gamma;t)=0$ for every closed curve $\gamma$ in $D$.
\item Every coefficient one-form of $\Om[g]$ is exact on $D$.
\end{enumerate}
When they hold, fix $z_0\in D$ and any $H_0\in\PP$. Then
\begin{equation}\label{eq:inverse-formula}
 H(z;t)=H_0(t)+\int_{z_0}^{z}\Om[g](t),\qquad F(z)=H(z;z)
\end{equation}
defines a holomorphic inverse. The integral is independent of path. Every
inverse occurs in this way, and any two differ by an element of $\PP$
evaluated at $z$.
\end{theorem}

\begin{proof}
If $g=\Fu F$, Theorem~\ref{thm:hermite} identifies the current with
$\mathrm dH_F$, so all coefficients are exact and all periods vanish.
Conversely, vanishing periods make the coefficientwise path integrals in
\eqref{eq:inverse-formula} path independent. Differentiation gives
$\mathrm dH=\Om[g]$. The standard elementary path-integral proof of exactness
applies componentwise; there are finitely many components.

It remains to prove that evaluating the resulting polynomial produces an
inverse. In central complex notation, \eqref{eq:complex-current} gives
\begin{align}
 \partial_zH(z;t)&=-\frac{\kappa_h}{2}(t-z)^{h-1}(t-\barz)^hW,
 \label{eq:H-z}\\
 \partial_{\barz}H(z;t)&=-\frac{\kappa_h}{2}(t-z)^h(t-\barz)^{h-1}\overline W.
 \label{eq:H-zbar}
\end{align}
For $F(z)=H(z;z)$ the $\barz$ derivative of the substituted variable $t=z$
is zero, and the right side of \eqref{eq:H-zbar} vanishes at $t=z$.
Thus $\partial_{\barz}F=0$.

We next establish the Hermite jets of $H$. The identity
$F(z)=H(z;z)$ is the case $j=0$ of
\begin{equation}\label{eq:reconstructed-jets}
 F^{(j)}(z)=\partial_t^jH(z;z),\qquad 0\le j<h.
\end{equation}
If $j<h-1$, differentiate this equality with respect to $z$. The total
derivative of its right side is
\[
 \partial_z\partial_t^jH(z;t)\big|_{t=z}
 +\partial_t^{j+1}H(z;z).
\]
The first term is zero: by \eqref{eq:H-z}, differentiating $j\le h-2$ times
in $t$ leaves at least one factor $t-z$. This proves the induction. For
$h=1$ there is no induction step and the value condition alone suffices.
The coefficients of $H$ are real, so the conjugate-node jets follow as well.
Uniqueness of Hermite interpolation now yields $H=H_F$.

Apply Theorem~\ref{thm:hermite} to the holomorphic stem just constructed:
\[
 \Om[\Fu F]=\mathrm dH_F=\mathrm dH=\Om[g].
\]
Lemma~\ref{lem:current-injective} gives $\Fu F=g$. Finally, if $F$ is any
inverse, its Hermite polynomial satisfies the same differential equation and
differs from the chosen $H$ by a constant element of $\PP$. This proves
both exhaustiveness and uniqueness modulo the kernel.
\end{proof}

\begin{corollary}[Local inversion and simply connected domains]
Every smooth axial target has a local holomorphic Fueter inverse. If $D$ is
simply connected, it has a global inverse on $D$.
\end{corollary}

\begin{proof}
On a sufficiently small disc, or on a simply connected planar domain, every
closed coefficient form has zero periods. Apply Theorem~\ref{thm:inverse}.
\end{proof}

There is no circular dependence on a pre-existing inverse theorem in this
argument: radial differentiation proves that the forward map lands in the
Vekua system; the Hermite identity is obtained from jets and two calibrations;
and the integral construction proves the inverse theorem directly.

\section{Polynomial monodromy, period realization, and covers}
\label{sec:cokernel}

\subsection{Continuation and its complete invariant}

On the universal cover of $D$, integrate the pullback of the current and use
\eqref{eq:inverse-formula}. Analytic continuation of this inverse around a
closed curve $\gamma$ adds the polynomial
\begin{equation}\label{eq:monodromy}
 F\longmapsto F+M_\gamma(z),\qquad
 M_\gamma(t)=\Per_g(\gamma;t).
\end{equation}
Indeed, continuation of the coefficient primitive $H$ adds the integral of
its differential around the curve, and subsequent evaluation at $t=z$ adds
exactly $M_\gamma(z)$.

Thus monodromy is additive and polynomial of degree below $2h$, with coefficients
in $E$. It is independent of which local inverse was chosen. The period map
\[
 \pi_1(D,z_0)\longrightarrow\PP,
 \qquad [\gamma]\longmapsto M_\gamma
\]
is a group homomorphism into an additive real vector group; it factors through
$H_1(D;\Z)$. Theorem~\ref{thm:inverse} says that it is a \emph{complete}
obstruction, not merely one invariant among several.

\subsection{Every finite period assignment occurs}

For the remainder of this subsection, assume that $D$ is bounded,
$\overline D\subset\C^+$, and its boundary consists of $b+1$ disjoint
$C^1$ Jordan curves. Select a point $p_j$ in each bounded complementary
component and select curves $\gamma_1,\ldots,\gamma_b$ in $D$ with
\[
 \operatorname{wind}(\gamma_i,p_j)=\delta_{ij}.
\]
Their homology classes form a basis of $H_1(D;\Z)$.

\begin{lemma}[Logarithmic representative]\label{lem:log-representative}
For $p\notin D$ and $M\in\PP$, define a local holomorphic stem by
\begin{equation}\label{eq:log-stem}
 F_{p,M}(z)=\frac{1}{2\pi\ii}M(z)\operatorname{Log}(z-p).
\end{equation}
Its Fueter image $G_{p,M}=\Fu F_{p,M}$ is single valued on $D$, independent
of branch, and belongs to $\AM(D;E)$. Moreover
\begin{equation}\label{eq:log-period}
 \Per_{G_{p,M}}(\gamma;t)=\operatorname{wind}(\gamma,p)M(t).
\end{equation}
\end{lemma}

\begin{proof}
Two logarithm branches differ by $2\pi\ii n$. The corresponding stems differ
by $nM(z)$, which lies in the real polynomial kernel of the Fueter map.
Their targets therefore agree, so the local targets patch to a smooth global
one. Continuation of the stem around $\gamma$ adds
$\operatorname{wind}(\gamma,p)M(z)$. The Hermite polynomial of a kernel
polynomial $M(z)$ is $M(t)$ itself. Integrating
$\mathrm dH_{F_{p,M}}=\Om[G_{p,M}]$ around the continued branch proves the
period formula, including its sign and normalization.
\end{proof}

\begin{remark}
The target $G_{p,M}$ need not be rational and need not be free of logarithmic
terms. Branch independence only says that changing $\operatorname{Log}$ by
$2\pi\ii n$ has zero Fueter image. Terms involving $\log|z-p|$ can remain.
No rationality assertion is required in any proof.
\end{remark}

\begin{theorem}[Explicit cokernel and splitting]\label{thm:cokernel}
Under the finite-connectivity hypotheses above, there is an exact sequence of
real vector spaces
\begin{equation}\label{eq:exact-sequence}
 0\longrightarrow\PP\longrightarrow\Hol(D;E_{\C})
 \xrightarrow{\ \Fu\ }\AM(D;E)
 \xrightarrow{\ \Per\ }\PP^b\longrightarrow0,
\end{equation}
where $\Per(g)=(\Per_g(\gamma_j;t))_{j=1}^b$. In particular,
\begin{equation}\label{eq:cokernel-dimension}
 \dim_{\R}\operatorname{coker}\Fu=2hdb.
\end{equation}
If $M_j=\Per_g(\gamma_j;t)$, then
\begin{equation}\label{eq:decomposition}
 g=\Fu F+\sum_{j=1}^bG_{p_j,M_j}
\end{equation}
for some global holomorphic stem $F$, unique modulo $\PP$.
\end{theorem}

\begin{proof}
Exactness at the first two nonzero terms follows from the kernel theorem.
Exactness at $\AM(D;E)$ is precisely Theorem~\ref{thm:inverse}, since the
chosen curves generate homology. Lemma~\ref{lem:log-representative} realizes
an arbitrary polynomial on one basis curve and zero on the others, proving
surjectivity of the last map. Subtract those representatives from $g$ and
apply the global inverse theorem to obtain \eqref{eq:decomposition}. Counting
the $2h$ coefficients in each of $b$ copies gives the dimension.
\end{proof}

For $h=1$ and $E=\HH$, \eqref{eq:cokernel-dimension} is $8b$, recovering the
quaternionic result. For $h=2$, there are four independent $E$-valued periods
per hole. The splitting depends on the chosen points and homology basis; the
period obstruction itself does not depend on a choice of local primitive.

\subsection{The smallest inverse cover and the finite-cover obstruction}

\begin{theorem}[Covering criterion]\label{thm:covers}
Let $g\in\AM(D;E)$, with no finite-connectivity assumption. A connected pointed
covering $\pi:\widetilde D\to D$ supports a single-valued inverse of the pulled-back
target, in the pulled-back local coordinate $z$, if and only if
\[
 \pi_*\pi_1(\widetilde D)\subseteq\ker\Per_g.
\]
The covering corresponding to $\ker\Per_g$ is the smallest such covering, in
the sense that every other connected inverse cover factors through it.
Its deck group is naturally $\operatorname{im}\Per_g$, regarded as an abstract
discrete group. If $\Per_g\ne0$, no finite-sheeted covering supports an inverse.
\end{theorem}

\begin{proof}
The integral construction and its necessity apply in the local coordinates
of a covering exactly as on $D$. The periods of the pulled-back current are
the original periods restricted to the image of its fundamental group. This
proves the criterion. The subgroup $\ker\Per_g$ is normal, so covering-space
classification gives the stated regular covering and deck group. Inclusion
of subgroups gives the factorization assertion.

For the last statement, let $\gamma$ have nonzero period $M$. If $H$ is any
finite-index subgroup of $\pi_1(D)$, some positive power $\gamma^N$ belongs to
$H$; for example, use the finite permutation action on cosets. Its period is
$NM\ne0$, because a real vector group has no nonzero torsion. Thus restriction
to $H$ cannot annihilate all periods.
\end{proof}

\begin{remark}[A topology caveat]
The image of the period homomorphism is a discrete group for purposes of the
deck action; it need not be a discrete subset of the Euclidean coefficient
space. For example, two polynomial periods can be irrationally proportional.
For finite $b$, the image is a finitely generated torsion-free abelian group
of rank at most $b$, even when its image in a real line is dense.
\end{remark}

\section{Coefficient formulas and the six-dimensional case}
\label{sec:algorithm}

The polynomial formulation is directly usable without expanding higher-order
inverses or choosing successive primitives. Write
\[
 C_j(x,y)=[t^j]\bigl((t-x)^2+y^2\bigr)^{h-1},
\]
with $C_j=0$ outside $0\le j\le2h-2$. Expansion of the three terms
$t^2$, $-2xt$, and $x^2+y^2$ gives
\begin{equation}\label{eq:C-coefficients}
 C_j=\sum_{\substack{a+b+c=h-1\\2a+b=j}}
 \frac{(h-1)!}{a!b!c!}(-2x)^b(x^2+y^2)^c.
\end{equation}
If $\Om[g](t)=\sum_{j=0}^{2h-1}\Omega_j[g]t^j$, then
\begin{equation}\label{eq:coefficient-forms}
\begin{split}
 \Omega_j[g]=\kappa_h\bigl[
 &\bigl((xC_j-C_{j-1})P+yC_jQ\bigr)\dd x\\
 &+\bigl(yC_jP+(C_{j-1}-xC_j)Q\bigr)\dd y\bigr].
\end{split}
\end{equation}
These forms use the values of $P,Q$, not their derivatives. The Vekua equations
are the proof that the forms are closed, not an additional step in evaluating
them along a curve.

\subsection{A constructive procedure}

For exact analytic input on a finitely connected domain, the procedure is:
compute the $2h$ coefficient forms, integrate them over a homology basis,
subtract logarithmic representatives of the resulting periods, integrate the
remaining exact forms along paths from a basepoint, and evaluate the resulting
polynomial at $t=z$. Formula~\eqref{eq:inverse-formula} proves the result,
including its kernel ambiguity.

Numerical data require an additional distinction. A small floating-point period
is not a proof that the period is zero. Certified claims require symbolic
integration, rigorous intervals together with a suitable exact zero criterion,
or another proof of vanishing. In particular, interval enclosures containing
zero by themselves do not certify exactness.

\subsection{Four explicit forms when the ambient dimension is six}

For $h=2$, $R_t=t^2-2xt+x^2+y^2$ and $\kappa_2=-1/16$. The complete four-form
system is
\begin{align}
 \Omega_0={}&-\frac{x^2+y^2}{16}
 \bigl[(xP+yQ)\dd x+(yP-xQ)\dd y\bigr],\label{eq:h2-0}\\
 \Omega_1={}&-\frac1{16}\bigl[
 ((-3x^2-y^2)P-2xyQ)\dd x\nonumber\\
 &\hspace{28mm}+(-2xyP+(3x^2+y^2)Q)\dd y\bigr],\label{eq:h2-1}\\
 \Omega_2={}&-\frac1{16}\bigl[(3xP+yQ)\dd x+(yP-3xQ)\dd y\bigr],
 \label{eq:h2-2}\\
 \Omega_3={}&\frac1{16}(P\dd x-Q\dd y).\label{eq:h2-3}
\end{align}
A global inverse exists exactly when all four forms have zero periods.
If $\mathrm dH_j=\Omega_j$, its stem is
\[
 F(z)=H_0(z)+H_1(z)z+H_2(z)z^2+H_3(z)z^3.
\]
Each $H_j$ is an $E$-valued real function; the displayed expression is a
central-complex evaluation, and its holomorphicity follows from the theorem,
not from separate holomorphicity of the $H_j$.

For the calibration target $P=64$, $Q=0$, which is
$\mathcal F_2(z^4)$, a coefficient primitive is
\[
 H(z;t)=t^4-\bigl((t-x)^2+y^2\bigr)^2.
\]
Although written as a difference of quartics, it is a cubic in $t$, and
$H(z;z)=z^4$. This is a simple exact check on all four forms simultaneously.

\section{Residue polynomials and a sharp growth filtration}
\label{sec:filtration}

Let $p=u+\ii v\in\C^+$ and choose $0<R<v$. Suppose $g\in\AM(D_R^*;E)$ on
\[
 D_R^*=\{z:0<|z-p|<R\}.
\]
Define its residue polynomial by
\begin{equation}\label{eq:residue-definition}
 M(t)=\int_{|z-p|=\rho}\Om[g](t),\qquad 0<\rho<R,
\end{equation}
where the circle is positively oriented. Closedness makes $M$ independent of
$\rho$. The excluded point $p$ corresponds physically to the conjugacy sphere
$\{u+vI:I\in S^{2h}\}$, not to a single point of the ambient space.

Put
\begin{equation}\label{eq:quadratic}
 \Delta_p(t)=(t-u)^2+v^2=(t-p)(t-\barp).
\end{equation}
For a vector-valued polynomial, vanishing of a jet means vanishing of all its
components. If $M\ne0$, define
\[
 \ord_pM=\min\{j\ge0:M^{(j)}(p)\ne0\}.
\]
Because its coefficients are real, a zero of order $r$ at $p$ is accompanied
by a zero of order $r$ at $\barp$. Consequently
\begin{equation}\label{eq:order-divisibility}
 \ord_pM\ge r\quad\Longleftrightarrow\quad
 M(t)\text{ is divisible in }E[t]\text{ by }\Delta_p(t)^r.
\end{equation}
This follows by scalar polynomial division in any real basis of $E$.
For nonzero $M\in\PP$, it implies $\ord_pM\le h-1$.

\subsection{A jet estimate extracted from the current}

\begin{lemma}[Residue jet bound]\label{lem:jet-bound}
For $0\le j\le h-1$, there is a constant $K_{h,p,j}$, independent of $g$
and of sufficiently small $\rho$, such that
\begin{equation}\label{eq:jet-bound}
 \norm{M^{(j)}(p)}
 \le K_{h,p,j}\rho^{h-j}
       \max_{|z-p|=\rho}\norm{W(z)}.
\end{equation}
\end{lemma}

\begin{proof}
Differentiate \eqref{eq:complex-current} $j$ times in $t$, set $t=p$, and
integrate over the circle. In the first term, the smallest possible power of
$p-z$ is $h-1-j$; in the second it is $h-j$. The other factor $p-\barz$
stays in a compact set bounded away from zero as $\rho\to0$, since
$p-\barp=2\ii v\ne0$. The arc element has length $\rho\dd\theta$.
Leibniz's rule, the triangle inequality, and the circle length $2\pi\rho$
therefore give \eqref{eq:jet-bound}. There are finitely many terms, so one
constant controls the whole sum.
\end{proof}

\begin{definition}[Growth-compatible residue space]
For an integer $0\le s\le h$, let $\mathcal R_s(p)$ be the set of polynomials
that occur as residues of at least one target satisfying
\[
 \norm{W(z)}=O(|z-p|^{-s})\quad(z\to p).
\]
This is a condition on the existence of a representative, not a claim that
every representative of the residue has that growth.
\end{definition}

\begin{theorem}[Exact residue filtration]\label{thm:filtration}
For $0\le s\le h$,
\begin{equation}\label{eq:filtration}
 \mathcal R_s(p)=\Delta_p(t)^{h-s}\PPl{2s},
 \qquad \dim_{\R}\mathcal R_s(p)=2sd.
\end{equation}
For a nonzero residue $M$, set $r=\ord_pM$ and $s=h-r$. There exists a target
with residue $M$ and growth $O(|z-p|^{-s})$, whereas no target with residue $M$
can have growth $O(|z-p|^{-\alpha})$ for any real $\alpha<s$.
\end{theorem}

\begin{proof}
Suppose first that $g$ has growth $O(\rho^{-s})$. For $j<h-s$,
Lemma~\ref{lem:jet-bound} gives
\[
 \norm{M^{(j)}(p)}=O(\rho^{h-j-s})\longrightarrow0.
\]
Since the left side is independent of $\rho$, it is zero. Thus
$\Delta_p^{h-s}$ divides $M$. The degree condition $\deg M<2h$ makes the
quotient have degree below $2s$. For $s=0$ this forces $M=0$.

For sufficiency, first take $M\ne0$ with $r=\ord_pM$ and $s_0=h-r$.
Use the logarithmic stem \eqref{eq:log-stem} on a punctured disc, with local
branches. Since $M$ vanishes to order $r$ at $p$, its derivatives after
multiplication by $\operatorname{Log}(z-p)$ satisfy, for $0\le j\le h$,
\[
 F_{p,M}^{(j)}(z)
 =O\!\left(1+|z-p|^{r-j}(1+|\log|z-p||)\right).
\]
At order $j=h>r$, differentiation of the lowest Taylor term gives a pole
of order $h-r=s_0$ without a logarithm in its leading term. All lower terms
have smaller order, with a possible logarithmic factor. The radial formula,
whose coefficients are smooth near $y=v>0$, therefore gives
$\norm{G_{p,M}}_{\mathrm{pair}}=O(|z-p|^{-s_0})$.
A fully normalized expansion is proved in Proposition~\ref{prop:canonical-asymptotic}.
If $M$ is divisible by $\Delta_p^{h-s}$, then $s_0\le s$, which proves
sufficiency for the claimed level. The zero residue is realized by the zero
target. Multiplication by a fixed nonzero scalar polynomial is injective on
$E[t]$, so the displayed space has dimension $2sd$.

Finally, for a nonzero residue use \eqref{eq:jet-bound} with $j=r$. Growth
$O(\rho^{-\alpha})$ for $\alpha<h-r$ would force its nonzero $r$th jet to
vanish. This is impossible.
\end{proof}

The successive quotients
$\mathcal R_s(p)/\mathcal R_{s-1}(p)$ all have dimension $2d$. In particular,
only a $2d$-dimensional subspace of the full $2hd$-dimensional residue space
can be represented by a simple-pole-growth target. This is the first
qualitative difference from the quaternionic case, where $h=1$ and every
nonzero obstruction is already on that first level.

\section{Optimal axial-pair singularity constants}
\label{sec:sharp}

The order filtration alone does not determine a leading coefficient. We now
show that its first nonzero jet determines an optimal one.

Fix $M\ne0$, let $r=\ord_pM$ and $s=h-r$, and define
\begin{equation}\label{eq:Aconstant}
 A_{h,p,M}=\frac{c_h(s-1)!}{2\pi v^h}\norm{M^{(r)}(p)}.
\end{equation}
For any target $g$ with residue $M$, put
\[
 B_g(\rho)=\max_{0\le\theta\le2\pi}
 \norm{W(p+\rho e^{\ii\theta})}.
\]

\begin{lemma}[Leading residue Fourier coefficient]\label{lem:fourier}
For sufficiently small $\rho$,
\begin{equation}\label{eq:fourier-extraction}
 M^{(r)}(p)=\gamma_{h,r,p}\rho^s
 \int_0^{2\pi}e^{\ii s\theta}W(p+\rho e^{\ii\theta})\dd\theta
 +\mathcal E(\rho),
\end{equation}
where
\begin{equation}\label{eq:gamma}
 \gamma_{h,r,p}=-\frac{\kappa_h}{2}
 \frac{(h-1)!}{(s-1)!}(-1)^{s-1}(2\ii v)^h\ii
\end{equation}
and, with a constant independent of the target,
\begin{equation}\label{eq:fourier-error}
 \norm{\mathcal E(\rho)}
 \le C\rho^{s+1}\int_0^{2\pi}
 \norm{W(p+\rho e^{\ii\theta})}\dd\theta.
\end{equation}
The modulus of the scalar in \eqref{eq:gamma} is
\begin{equation}\label{eq:gamma-modulus}
 b_{h,r,p}=|\gamma_{h,r,p}|=\frac{v^h}{c_h(s-1)!}.
\end{equation}
\end{lemma}

\begin{proof}
In the first term of \eqref{eq:complex-current}, differentiate $r$ times in
$t$ and set $t=p$. The unique lowest-order contribution in $p-z$ occurs when
all $r$ derivatives hit $(t-z)^{h-1}$. It is
\[
 -\frac{\kappa_h}{2}\frac{(h-1)!}{(s-1)!}
 (p-z)^{s-1}(p-\barz)^hW\dd z.
\]
On the circle, $p-z=-\rho e^{\ii\theta}$,
$\dd z=\ii\rho e^{\ii\theta}\dd\theta$, and
$p-\barz=2\ii v-\rho e^{-\ii\theta}$. Replacing its $h$th power by
$(2\ii v)^h$ produces the leading term in \eqref{eq:fourier-extraction};
the replacement error has one more power of $\rho$.

Every other distribution of derivatives in the first term leaves at least
$s$ factors of $p-z$. The entire second term of the complex current also
has at least $s$ such factors after differentiation. The additional arc
factor $\rho$ gives the common error estimate \eqref{eq:fourier-error}.
The bounded remaining factors and finitely many coefficients are absorbed
into $C$. Finally, inserting the value of $\kappa_h$ into
$|\gamma_{h,r,p}|$ and cancelling factorials gives \eqref{eq:gamma-modulus}.
\end{proof}

\Needspace{15\baselineskip}
\begin{theorem}[Sharp size and circle-energy bounds]\label{thm:sharp-size}
Every smooth axial target with nonzero residue $M$ satisfies
\begin{equation}\label{eq:sharp-sup}
 \liminf_{\rho\downarrow0}\rho^s B_g(\rho)\ge A_{h,p,M}.
\end{equation}
More precisely, for small $\rho$,
\begin{equation}\label{eq:circle-L2}
 \int_0^{2\pi}\norm{W(p+\rho e^{\ii\theta})}^2\dd\theta
 \ge 2\pi A_{h,p,M}^{\,2}\rho^{-2s}(1-O(\rho)).
\end{equation}
No upper growth assumption on $g$ is needed for either conclusion.
Both constants are attained asymptotically by $G_{p,M}$ from
Lemma~\ref{lem:log-representative}.
\end{theorem}

\begin{proof}[Proof of the lower bounds]
Lemma~\ref{lem:fourier} and the triangle inequality imply
\begin{equation}\label{eq:circle-L1}
 \norm{M^{(r)}(p)}
 \le(b_{h,r,p}+C\rho)\rho^s
 \int_0^{2\pi}\norm{W(p+\rho e^{\ii\theta})}\dd\theta.
\end{equation}
The integral is at most $2\pi B_g(\rho)$, which gives
\eqref{eq:sharp-sup} because
$\norm{M^{(r)}(p)}/(2\pi b_{h,r,p})=A_{h,p,M}$.
Alternatively, apply Cauchy--Schwarz to the integral in
\eqref{eq:circle-L1}, square, and rearrange:
\[
 \int_0^{2\pi}\norm W^2\dd\theta
 \ge\frac{\norm{M^{(r)}(p)}^2}
 {2\pi(b_{h,r,p}+C\rho)^2\rho^{2s}}.
\]
Expansion of the scalar denominator proves \eqref{eq:circle-L2}.
These estimates remain valid for a target growing faster than any power;
they used only integrals on individual compact circles.
\end{proof}

\begin{proposition}[Canonical leading asymptotic]\label{prop:canonical-asymptotic}
The target of the logarithmic stem \eqref{eq:log-stem} has axial pair
$W_{p,M}$ satisfying
\begin{equation}\label{eq:canonical-leading}
\begin{split}
 W_{p,M}(z)={}&
 \frac{c_h\ii^{h-1}(-1)^{s-1}(s-1)!}{2\pi v^h}
 M^{(r)}(p)(z-p)^{-s}\\
 &+O\!\left(|z-p|^{-s+1}(1+|\log|z-p||)\right).
\end{split}
\end{equation}
The remainder estimate is uniform in the angular variable. Hence
\begin{align}
 \lim_{\rho\downarrow0}\rho^sB_{G_{p,M}}(\rho)&=A_{h,p,M},
 \label{eq:canonical-sup}\\
 \lim_{\rho\downarrow0}\rho^{2s}\int_0^{2\pi}
 \norm{W_{p,M}(p+\rho e^{\ii\theta})}^2\dd\theta
 &=2\pi A_{h,p,M}^2.\label{eq:canonical-circle}
\end{align}
\end{proposition}

\begin{proof}
Expand the real-coefficient polynomial at $p$:
\[
 M(z)=\frac{M^{(r)}(p)}{r!}(z-p)^r+O((z-p)^{r+1}).
\]
For $h=r+s$ and $s\ge1$,
\[
 \frac{\mathrm d^h}{\mathrm dz^h}\bigl((z-p)^r\operatorname{Log}(z-p)\bigr)
 =r!(-1)^{s-1}(s-1)!(z-p)^{-s}.
\]
Thus the $h$th derivative of the logarithmic stem is
\begin{equation}\label{eq:log-high-derivative}
\begin{split}
 F_{p,M}^{(h)}(z)={}&
 \frac{(-1)^{s-1}(s-1)!}{2\pi\ii}
 M^{(r)}(p)(z-p)^{-s}\\
 &+O\!\left(|z-p|^{-s+1}(1+|\log|z-p||)\right).
\end{split}
\end{equation}
Use a branch with argument in an interval of length $2\pi$ for this estimate.
The resulting target is branch independent, so the bounds extend uniformly
around the circle.

Both radial operators have leading term $y^{-h}\partial_y^h$ after $h$
iterations. All other terms have $y$-smooth coefficients and derivatives of
order at most $h-1$. Since a holomorphic stem satisfies
$\partial_y^hF=\ii^hF^{(h)}$, the pair form of \eqref{eq:radial} gives
\[
 W_{p,M}=c_h\ii^hy^{-h}F_{p,M}^{(h)}
   +\text{terms involving derivatives of order at most }h-1.
\]
The lower-order terms obey the remainder estimate in
\eqref{eq:canonical-leading}. Replacing $y^{-h}$ by $v^{-h}$ also changes
only that remainder, because $y-v=O(|z-p|)$. Formula
\eqref{eq:log-high-derivative} now yields \eqref{eq:canonical-leading}.
The scalar angular factor $e^{-\ii s\theta}$ has modulus one and acts
isometrically on $E_{\C}$. The norm of the leading coefficient is exactly
$A_{h,p,M}$; the remainder divided by $\rho^{-s}$ tends uniformly to zero.
Both limits follow.
\end{proof}

\subsection{Normalization against the quaternionic case}

If $h=1$ and $M(t)=ta+b$, then $r=0$, $s=1$, and
\[
 A_{1,p,M}=\frac1{\pi v}\norm{pa+b}
 =\frac1{\pi v}\sqrt{\norm{ua+b}^2+v^2\norm a^2}.
\]
Here the first norm is in the central complexification of $E$. If $E=\HH$,
$p=\ii$, $a=1$, and $b$ is a unit quaternionic imaginary vector orthogonal to
$1$, this constant is $\sqrt2/\pi$. A whole-conjugacy-sphere supremum is a
different norm and can have a different sharp constant. This explicitly
preserves the distinction emphasized by the repository's quaternionic report
\cite{repo}.

\section{Sharp physical energy and the integrability obstruction}
\label{sec:energy}

Let
\[
 \sig=\operatorname{area}(S^{2h})
       =\frac{2\pi^{h+1/2}}{\Gamma(h+1/2)}.
\]
For $0<\varepsilon<R<v$, define the tube around the punctured sphere by
\[
 T_{\varepsilon,R}(p)
 =\{x+yI:\varepsilon<|x+\ii y-p|<R,\ I\in S^{2h}\}.
\]
Its Euclidean volume element is
$\dd q=y^{2h}\dd x\dd y\dd\sigma(I)$. The physical energy is
\begin{equation}\label{eq:physical-energy}
 \mathcal E_g(\varepsilon,R)
 =\int_{T_{\varepsilon,R}(p)}\norm{g(q)}^2\dd q.
\end{equation}

\begin{lemma}[Exact spherical averaging]\label{lem:spherical-average}
For all $P,Q\in E$,
\begin{equation}\label{eq:spherical-average}
 \int_{S^{2h}}\norm{P+IQ}^2\dd\sigma(I)
 =\sig\bigl(\norm P^2+\norm Q^2\bigr).
\end{equation}
Consequently
\begin{equation}\label{eq:pair-energy}
 \mathcal E_g(\varepsilon,R)
 =\sig\int_{\varepsilon<|z-p|<R}y^{2h}\norm{W(z)}^2\dd x\dd y.
\end{equation}
\end{lemma}

\begin{proof}
The identity $\norm{IQ}=\norm Q$ follows from the orthogonal module structure.
The cross term $2\langle P,IQ\rangle$ changes sign under $I\mapsto-I$ and
therefore has zero spherical integral. The volume formula gives the second
assertion.
\end{proof}

\Needspace{17\baselineskip}
\begin{theorem}[Optimal energy divergence]\label{thm:energy}
Let $g$ have nonzero residue $M$ at $p$, with $r=\ord_pM$ and $s=h-r$.
Fix $R$ with $0<R<v$. If $s=1$, then
\begin{equation}\label{eq:energy-log}
 \liminf_{\varepsilon\downarrow0}
 \frac{\mathcal E_g(\varepsilon,R)}{\log(1/\varepsilon)}
 \ge \frac{\sig c_h^2}{2\pi}\norm{M^{(h-1)}(p)}^2.
\end{equation}
If $s\ge2$, then
\begin{equation}\label{eq:energy-power}
 \liminf_{\varepsilon\downarrow0}
 \varepsilon^{2s-2}\mathcal E_g(\varepsilon,R)
 \ge\frac{\sig c_h^2((s-1)!)^2}{4\pi(s-1)}
       \norm{M^{(r)}(p)}^2.
\end{equation}
The logarithmic representatives $G_{p,M}$ attain equality as limits in both
formulas. In particular, nonzero polynomial residue forces infinite physical
$L^2$ energy near the conjugacy sphere.
\end{theorem}

\begin{proof}
Use polar coordinates $z=p+\rho e^{\ii\theta}$ in
\eqref{eq:pair-energy}. Uniformly in $\theta$,
\[
 y^{2h}=v^{2h}(1+O(\rho)).
\]
The circle lower bound \eqref{eq:circle-L2} therefore gives, on a sufficiently
small fixed radial interval,
\[
 \mathcal E_g(\varepsilon,R)
 \ge2\pi\sig v^{2h} A_{h,p,M}^2
 \int_\varepsilon^{R_0}\rho^{1-2s}(1-O(\rho))\dd\rho.
\]
For $s=1$, the main integral is $\log(R_0/\varepsilon)$, and its error is
bounded. For $s\ge2$, its main divergent term is
$\varepsilon^{2-2s}/(2s-2)$, while the error is lower order. Substitution of
\eqref{eq:Aconstant} cancels the factor $v^{2h}$ and gives the stated
constants. The discarded outer annulus has nonnegative energy, so it does
not weaken these lower bounds.

For $G_{p,M}$, Proposition~\ref{prop:canonical-asymptotic} gives the matching
uniform leading integrand. The relative error is
$O(\rho(1+|\log\rho|))$, whose integral is lower order in both regimes.
The energy on a fixed outer annulus is finite because the target is smooth
there. Hence equality holds as a full limit.
\end{proof}

\subsection{An exact local \texorpdfstring{$L^p$}{Lp} classification of residue classes}

The same argument has a consequence beyond the quadratic norm. To avoid
confusing an exponent with the point $p$, write the integrability exponent
as $\alpha$.

\begin{corollary}[$L^\alpha$-compatible residues, $\alpha\ge1$]
\label{cor:Lp}
Fix a nonzero residue $M$, with $s=h-\ord_pM$. There exists a representative
with finite physical $L^\alpha$ norm near the punctured sphere if and only if
\begin{equation}\label{eq:Lp-threshold}
 \alpha s<2.
\end{equation}
Equivalently, the spaces of residues admitting an $L^\alpha$ representative
are
\begin{equation}\label{eq:Lp-spaces}
 \begin{cases}
  \Delta_p^{h-1}\PPl{2},&1\le\alpha<2,\\
  \{0\},&\alpha\ge2.
 \end{cases}
\end{equation}
This classifies attainable residue classes, not all singular functions with
a given residue.
\end{corollary}

\begin{proof}
For a fixed pair $P,Q$, averaging $P+IQ$ over $I$ recovers $P$, and averaging
$-I(P+IQ)$ recovers $Q$. By the triangle inequality and the isometry of $I$,
\[
 \frac1{\sig}\int\norm{P+IQ}\dd\sigma(I)
 \ge\max(\norm P,\norm Q)\ge2^{-1/2}\norm W.
\]
Jensen's inequality for $\alpha\ge1$ then gives a lower bound
$\sig 2^{-\alpha/2}\norm W^\alpha$ for the integral of
$\norm{P+IQ}^\alpha$ over the sphere. The elementary upper bound
$\norm{P+IQ}\le\sqrt2\norm W$ gives a matching equivalence up to constants.
Thus physical $L^\alpha$ integrability is equivalent to integrability of
$y^{2h}\norm W^\alpha$ in the planar variables.

Equation~\eqref{eq:circle-L1} and H\"older's inequality imply
\[
 \int_0^{2\pi}\norm W^\alpha\dd\theta
 \ge2\pi A_{h,p,M}^{\alpha}\rho^{-\alpha s}(1-O(\rho)).
\]
Its radial integral diverges when $\alpha s\ge2$, including equality. On the
other hand, the canonical target has $\norm W\sim A_{h,p,M}\rho^{-s}$
uniformly, so it is integrable precisely when $\alpha s<2$.
For $\alpha\ge1$ and integer $s\ge1$, this condition means $s=1$ and
$\alpha<2$. Apply the residue filtration. The zero residue is realized by
the zero target for every exponent.
\end{proof}

\subsection{What zero residue does not imply}

Finite physical $L^2$ energy implies that the period polynomial is zero, so
a single-valued holomorphic Fueter inverse exists on the punctured disc.
Neither implication asserts that the target or the inverse extends through
the puncture. For example, the single-valued stem
\[
 F(z)=\frac{a}{z-p},\qquad a\in E\setminus\{0\},
\]
has zero monodromy and hence zero polynomial residue, but its Fueter image
has leading axial-pair order $|z-p|^{-h-1}$ by the radial formula. It is
singular and has infinite energy. A removability theorem needs additional
analytic hypotheses; the period obstruction alone is topological.

\section{Worked singularity examples in dimension six}
\label{sec:examples}

Take $h=2$, $p=\ii$, and a unit real coefficient $a\in E$. Here $c_2=8$,
$\kappa_2=-1/16$, and $\sigma_4=8\pi^2/3$. The residue space is the space
of cubic polynomials with coefficients in $E$.

\subsection{A residue that forces a double-pole scale}

Let $M(t)=a$. Its first nonzero jet is $M(p)=a$, so $r=0$ and $s=2$.
The canonical stem is
\[
 F(z)=\frac{a}{2\pi\ii}\operatorname{Log}(z-\ii).
\]
Every target with this residue satisfies
\[
 \liminf_{\rho\downarrow0}\rho^2B_g(\rho)\ge\frac4\pi.
\]
The logarithmic representative attains equality. Its optimal physical energy
constant, and the lower bound for every representative, are
\[
 \liminf_{\varepsilon\downarrow0}\varepsilon^2\mathcal E_g(\varepsilon,R)
 \ge\frac{128\pi}{3}.
\]
No simple-pole-growth representative exists. In fact no representative has
finite physical $L^1$ norm near the sphere, by Corollary~\ref{cor:Lp}.

\subsection{A quadratic factor lowers the forced order}

Let $M(t)=(t^2+1)a$. Then $M(p)=0$, $M'(p)=2\ii a$, so $r=1$ and $s=1$.
A canonical stem is
\[
 F(z)=\frac{(z^2+1)a}{2\pi\ii}\operatorname{Log}(z-\ii).
\]
The minimum axial-pair leading size is now
\[
 \liminf_{\rho\downarrow0}\rho B_g(\rho)\ge\frac8\pi,
\]
and the optimal energy law is
\[
 \liminf_{\varepsilon\downarrow0}
 \frac{\mathcal E_g(\varepsilon,R)}{\log(1/\varepsilon)}
 \ge\frac{1024\pi}{3}.
\]
Both are attained by the canonical target. It belongs locally to every
physical $L^\alpha$ space with $1\le\alpha<2$, but not to $L^2$.

The full simple-pole-compatible residue subspace in this dimension is
\[
 (t^2+1)\{a_0+ta_1:a_0,a_1\in E\},
\]
of dimension $2d$ inside a $4d$-dimensional total residue space. Thus the
quadratic divisibility condition is not a cosmetic rewriting of a norm
estimate: it removes half the independent obstruction directions at the
simple-pole level.

\subsection{General dimensions and higher levels}

At $p=\ii$, the family
\[
 M_s(t)=(t^2+1)^{h-s}a,\qquad 1\le s\le h,
\]
realizes every possible least growth order. Its first nonzero jet is
\[
 M_s^{(h-s)}(\ii)=(h-s)!(2\ii)^{h-s}a.
\]
Therefore
\[
 A_{h,\ii,M_s}=
 \frac{2^hh!(s-1)!(h-s)!2^{h-s}}{2\pi}\norm a.
\]
These are the families used in the numerical asymptotic diagnostics shipped
with the article. Their purpose is to test constants at all levels, not to
supply the general proof of the filtration.

\section{Verification, proof dependencies, and scope of the research claim}
\label{sec:audit}

\subsection{The all-order proofs versus finite diagnostics}

The code in \texttt{code/verify.py} was executed for this package. It uses
SymPy for exact symbolic identities and mpmath at 50 decimal digits for
numerical diagnostics. The complete retained output is in
\texttt{data/verification.json}, with a readable run log alongside it.

The exact checks establish closedness after substituting the Vekua system for
$1\le h\le8$; the Hermite differential identity for 60 real-monomial cases
($1\le h\le5$, $0\le k\le2h+5$); 13 additional imaginary-monomial cases;
and calibration and leading-residue phase identities through $h=8$.
The imaginary tests are especially useful because they detect an accidental
replacement of the real polynomial kernel by a complex polynomial kernel.

The numerical tests integrate every coefficient of the period current for a
nontrivial logarithmic representative in each of $h=1,2,3,4$, using 384 equally
spaced quadrature nodes. The largest observed absolute error over these tests
was below $4\times10^{-48}$. The tests also sample all ten pairs
$1\le s\le h\le4$ at four radii from $10^{-2}$ to $10^{-5}$. At the smallest
radius, the sampled normalized maximum in every case differed from the
predicted limit by less than $3.7\times10^{-4}$. These are observed numerical
errors, not rigorous enclosures of quadrature errors or certified bounds on
unsampled extrema.

None of those finite ranges is used to prove an unbounded assertion in the
article. The general Hermite identity follows from the exact jet-spanning
argument in Theorem~\ref{thm:hermite}; its reduction to $2h+2$ monomials is
valid for every symbolic $h$. The sharp constants follow from the explicit
Fourier extraction and logarithmic asymptotic. This is a logical distinction,
not simply a difference in precision.

\subsection{A dependency map}

The argument has two main chains:
\[
 \text{radial calculus}
 \ \Longrightarrow\ \text{two calibrations}
 \ \Longrightarrow\ \text{Hermite identity}
 \ \Longrightarrow\ \text{global inverse},
\]
and
\[
 \text{polynomial current}
 \ \Longrightarrow\ \text{residue jet extraction}
 \ \Longrightarrow\ \text{growth and energy lower bounds}.
\]
The logarithmic representative connects them: the polynomial kernel guarantees
its target is single valued, the Hermite identity computes its residue, and
ordinary differentiation computes its leading singularity. This supplies
attainment, rather than only a universal lower estimate.

\begin{center}
\begin{tabular}{p{0.28\textwidth}p{0.60\textwidth}}
\toprule
Conclusion & Essential ingredients \\
\midrule
Complete period criterion & Closedness, Hermite jets, current injectivity \\
Finite cokernel and splitting & Kernel theorem, logarithm monodromy, planar homology \\
No finite inverse cover & Additivity of periods, torsion-freeness \\
Residue filtration & Jet estimate, real conjugate roots, logarithmic representative \\
Sharp size constant & Leading Fourier coefficient and its exact normalization \\
Sharp physical energy & Circle estimate, spherical average, radial integration \\
$L^\alpha$ residue classes & Circle $L^1$ estimate, Jensen/H\"older, attainment \\
\bottomrule
\end{tabular}
\end{center}

\subsection{Repository and literature boundaries}

The repository source inspected for this extension was its merged
quaternionic-inversion README, in the tree identified by commit
\texttt{fbba58593dc0622aa914972896150d4848f935b5}; the report's path is recorded
in the bibliography and the package README. This was a targeted inspection,
not a proof-level audit of every source manuscript or every file in ProveIt.
All repository formulas used as a starting point have been rederived here
in the $h=1$ specialization.

Colombo, Pe\~na Pe\~na, Sabadini, and Sommen give integral inverses through
radial operators and coefficient ordinary differential equations
\cite{colombo2014}; their Theorem~3 is formulated on a rectangular radial
region bounded away from the axis. Colombo, Sabadini, and Sommen develop an
integral inverse using spherical monogenics \cite{colombo2013}. De Martino,
Diki, and Guzm\'an Ad\'an study Fueter--Sce images, generalized
Cauchy--Kovalevskaya extension, and the polynomial kernel
\cite{demartino2023}. Those works provide the primary theoretical context;
this article does not rebrand the forward mapping theorem, local inversion,
or the polynomial kernel as discoveries.

For the quaternionic global topology comparison, the explicit simply connected
hypothesis in Perotti's Theorem~2 is the relevant statement
\cite{perotti2024}. It is compatible with our result. We do not infer a false
unconditional surjectivity theorem from an abstract, introductory sentence,
or a theorem quoted without its domain assumptions.

The main proposed contribution should therefore be judged at the level of
its actual formulas and implications: the target-only current of degree
$2h-1$, the Hermite reconstruction, the explicit monodromy realization, and
the sharp hierarchy of residue-forced singularities. Whether these exact
results appear elsewhere under a different formalism remains a literature
question. The proofs establish the statements under the given hypotheses;
they do not establish historical priority.

\subsection{Formalization opportunities and present status}

No Lean source is included or claimed verified. A future formalization can be
split at clear interfaces. The polynomial portion consists of conjugate-node
Hermite interpolation, two remainders, and a finite-jet linearity theorem.
The differential portion consists of the radial commutators, Cauchy--Riemann
identities, and the coefficient-current calculation. The topological portion
is coefficientwise exactness versus periods. Finally, the quantitative portion
requires normed finite-dimensional complexification, a circle integral
inequality, and explicit improper-integral asymptotics.

This division would make a formal audit informative: one could certify the
algebraic construction before formalizing covering spaces or singular
integrals. Existing formal theorems elsewhere in the repository confer no
formal status on the new results in this article.

\section{Further research questions and topics}
\label{sec:future}

The following questions are continuations of the proved results, not assumptions
used in their proofs. Several concern quantitative or constructive refinements
of otherwise familiar general principles.

\begin{enumerate}[label=\textbf{\arabic*.},leftmargin=2em,itemsep=7pt]
\item \textbf{Crossing the real axis.}
The interpolation nodes $z$ and $\barz$ coalesce at $y=0$. Find a regularized
current and an exact global criterion on axial domains meeting the real axis,
with the required even/odd stem compatibility included explicitly. Determine
which period coefficients become redundant or constrained when cycles can be
moved through the axis. The present formulas do not settle this degeneration.

\item \textbf{Higher angular monogenic degree.}
For a fixed spherical monogenic of degree $\ell$, the radial Vekua coefficient
is shifted from $2h$ to $2(h+\ell)$. Determine the correct target-only current,
normalization, and real coefficient module for the corresponding weighted
Fueter map. Establish which of the residue filtration and sharp physical
energy constants survive after angular orthogonality and multiplicity are
accounted for. Merely replacing $h$ in one display does not by itself verify
the forward-map normalization or the energy measure.

\item \textbf{Global least-energy representatives.}
For a multiply connected smooth domain and fixed period polynomials, construct
the least-energy target under natural boundary conditions. Determine the
finite-dimensional quadratic form on the period space and its relation to a
capacity or period matrix. Near shrinking holes, prove that its diagonal
terms have the sharp divergent constants of Theorem~\ref{thm:energy} and
compute the finite interaction terms.

\item \textbf{Stable inverse operators.}
After subtracting a fixed period splitting and fixing a polynomial gauge,
prove norm estimates for the inverse on Sobolev or H\"older spaces. Identify
the dependence on boundary regularity, the distance to the real axis, and
neck degeneration in the planar domain. An explicit formula alone does not
provide a uniformly bounded inverse in a chosen norm.

\item \textbf{Infinite connectivity with growth control.}
The zero-period criterion already holds on every connected $D\subset\C^+$.
For infinitely many holes, characterize which polynomial period assignments
admit representatives satisfying specified growth or summability conditions.
Develop convergent analogues of the finite logarithmic splitting, including
estimates near accumulating boundary components. Unrestricted sheaf-theoretic
existence and constructive growth-controlled realization should be kept
separate.

\item \textbf{Collision and real-axis limits of residues.}
Study families of punctures that approach one another or approach the real
axis. Determine how the quadratic factors $\Delta_p$ and their filtrations
combine, and derive a uniform asymptotic transition between distinct spherical
singularities and the limiting singular set. The fixed-$v>0$ estimates here
are not uniform as $v\to0$.

\item \textbf{Nonaxial targets.}
Expand a general monogenic function into angular modes and ask when modewise
Fueter inverses can be recombined with convergence and compatible global
monodromy. A useful theorem would give a common analytic class and a complete
obstruction space, not just formal independent inverses of each mode.

\item \textbf{Fractional Fueter maps.}
In dimensions where the natural Fueter map is expressed through a fractional
Laplacian rather than an integer power, determine whether an analogous
finite-dimensional obstruction survives. The proof here relies on finite
jets and a finite polynomial kernel; a nonlocal map may require a genuinely
different invariant rather than more coefficient forms.

\item \textbf{Endpoint spaces and stronger removability.}
The strong $L^\alpha$ residue classification for $\alpha\ge1$ is settled here.
Determine optimal weak-$L^\alpha$, Lorentz, and weighted analogues, and identify
additional hypotheses that turn zero residue into removability. Separate
conditions on the target from conditions on a selected holomorphic inverse,
since a zero-residue target can still have arbitrarily severe singularities.

\item \textbf{Certified symbolic and formal inversion.}
Implement the coefficient-current construction, homology-basis integration,
and polynomial gauge choice with exact certificates. A first proof-assistant
target is the universal Hermite differential identity, followed by the global
criterion and the sharp coefficient extraction. For numerical input, develop
an error model that distinguishes small periods from provably zero periods.
\end{enumerate}

\section{Conclusion}

The higher-dimensional inverse Fueter problem on upper-half-plane product
domains is controlled by a finite polynomial object. Its coefficients are
closed one-forms computed directly from the target. Its primitive is the
Hermite polynomial of any local holomorphic inverse, and evaluating that
polynomial at the moving node reconstructs the inverse. This proves
completeness of the periods and realizes arbitrary finite period data by
explicit logarithmic stems.

The same current carries quantitative information. The first nonzero
conjugate-node jet of a residue polynomial fixes the least possible growth
order; its norm fixes an optimal leading size. Spherical averaging then gives
sharp physical energy constants. These results extend the quaternionic
pair-of-periods picture to a $2h$-coefficient obstruction and reveal a
$h$-level residue hierarchy that has no nontrivial intermediate levels when
$h=1$.

The remaining limits are explicit: this is axial analysis on domains avoiding
the real axis; the finite-dimensional cokernel realization was made
constructive for finite connectivity; the sharp norm is the axial-pair norm;
and zero residue is not a removability theorem. Within these boundaries, the
article gives proofs rather than conjectural extrapolations.

\appendix
\section{An explicit formula for the Hermite interpolant}
\label{app:hermite}

The interpolation polynomial used in the proof can be written without solving
a linear system. For fixed $z\in D$, regard $w$ as an independent complex
variable, and let $F^*(w)=\overline{F(\overline w)}$ on a neighborhood of
$\barz$. Then
\begin{equation}\label{eq:explicit-hermite}
\begin{split}
 H_F(z;t)={}&(t-\barz)^h\sum_{j=0}^{h-1}
 \frac{(t-z)^j}{j!}
 \left.\frac{\mathrm d^j}{\mathrm dw^j}
       \left(\frac{F(w)}{(w-\barz)^h}\right)\right|_{w=z}\\
 &+(t-z)^h\sum_{j=0}^{h-1}
 \frac{(t-\barz)^j}{j!}
 \left.\frac{\mathrm d^j}{\mathrm dw^j}
       \left(\frac{F^*(w)}{(w-z)^h}\right)\right|_{w=\barz}.
\end{split}
\end{equation}
The first term matches the required Taylor jet at $z$ and vanishes to order
$h$ at $\barz$; the second does the opposite. Each has degree at most $2h-1$.
Their sum is invariant under conjugation of coefficients. This proves the
formula and also exhibits its dependence on only the first $h$ stem jets.

For a real-coefficient polynomial $F(w)$, $H_F$ is simply the remainder upon
division of $F(t)$ by $R_t^h$. That observation is used in the exact
real-monomial tests. For imaginary-coefficient stems, ordinary division of
$F(t)$ by $R_t^h$ need not give a real-coefficient interpolant, so the tests
instead use \eqref{eq:explicit-hermite}.

\section{Explicit radial coefficient formulas}
\label{app:radial}

For computational use, repeated radial differentiation can be written as
\begin{align}
 \Rop^h
 &=\sum_{j=1}^h
 \frac{(-1)^{h+j}(2h-j-1)!}{2^{h-j}(h-j)!(j-1)!}
 y^{j-2h}\partial_y^j,\label{eq:R-explicit}\\
 \Sop^h
 &=\sum_{j=0}^h
 \frac{(-1)^{h+j}(2h-j)!}{2^{h-j}(h-j)!j!}
 y^{j-2h}\partial_y^j.\label{eq:S-explicit}
\end{align}
For $h=1$, these are their defining formulas. Inductively, applying
$y^{-1}\partial_y$ to a term $a_{h,j}y^{j-2h}\partial_y^j$ gives the
coefficient recurrence
\[
 a_{h+1,j}=(j-2h)a_{h,j}+a_{h,j-1}.
\]
Applying $\partial_y(y^{-1}\,\cdot\,)$ gives
\[
 b_{h+1,j}=(j-2h-1)b_{h,j}+b_{h,j-1}.
\]
Missing coefficients are zero. Substitution of the factorial expressions
verifies these recurrences and proves both formulas.

For a logarithmic representative, derivatives can be computed exactly from
\begin{equation}\label{eq:log-jets-algorithm}
 F_{p,M}^{(j)}(z)=\frac1{2\pi\ii}
 \sum_{\ell=0}^j\binom j\ell M^{(j-\ell)}(z)L_\ell(z),
\end{equation}
where
\[
 L_0(z)=\operatorname{Log}(z-p),\qquad
 L_\ell(z)=(-1)^{\ell-1}(\ell-1)!(z-p)^{-\ell}\quad(\ell\ge1).
\]
For holomorphic stems $\partial_y^jF=\ii^jF^{(j)}$, so taking real and imaginary
parts in \eqref{eq:R-explicit}--\eqref{eq:S-explicit} avoids numerical
finite-difference differentiation. This is the method used in the period and
asymptotic diagnostics.

\section{Reproduction and package contents}

The package contains the complete \texttt{article.tex} and compiled
\texttt{article.pdf}, the verification program, its recorded JSON output and
run log, a short proof audit, and build instructions. The bibliography is
embedded in the source, so no separate bibliography processor is required.

Rebuild the article using
\begin{lstlisting}
sh build.sh
\end{lstlisting}
and rerun the complete diagnostic suite using
\begin{lstlisting}
python code/verify.py
\end{lstlisting}
The verification script requires SymPy and mpmath. The recorded run used
SymPy 1.14.0 and mpmath 1.3.0. A \texttt{--quick} switch reduces only the
range of real-monomial symbolic checks; it does not change the statements
proved in the article. Running the script replaces its JSON output with the
newly observed results. There are no hidden certificate files and no
probabilistic claims of exact nonvanishing.

\begin{thebibliography}{9}
\raggedright

\bibitem{repo}
\emph{ProveIt} project, \emph{Slice regularity and the global inverse Fueter
problem}, repository report README, September 2026.
Path: \texttt{SetTheory/Cardinals/docs/reports/quaternionic-analysis/}\\
\texttt{slice-regularity-and-fueter-inversion/README.md}.
Inspected repository snapshot:
\texttt{fbba58593dc0622aa914972896150d4848f935b5}.
\href{https://github.com/VladimirReshetnikov/ProveIt/tree/fbba58593dc0622aa914972896150d4848f935b5/SetTheory/Cardinals/docs/reports/quaternionic-analysis/slice-regularity-and-fueter-inversion}{Pinned repository report}.
Unrefereed research materials; the present article rederives the formulas used.

\bibitem{colombo2014}
F.~Colombo, D.~Pe\~na Pe\~na, I.~Sabadini, and F.~Sommen,
\emph{A new integral formula for the inverse Fueter mapping theorem},
Journal of Mathematical Analysis and Applications \textbf{417} (2014),
112--122.
\href{https://doi.org/10.1016/j.jmaa.2014.03.016}{doi:10.1016/j.jmaa.2014.03.016}.
Preprint: \href{https://arxiv.org/abs/1302.0685}{arXiv:1302.0685}.

\bibitem{colombo2013}
F.~Colombo, I.~Sabadini, and F.~Sommen,
\emph{The inverse Fueter mapping theorem in integral form using spherical
monogenics}, Israel Journal of Mathematics \textbf{194} (2013), 485--505.
\href{https://doi.org/10.1007/s11856-012-0090-4}{doi:10.1007/s11856-012-0090-4}.

\bibitem{demartino2023}
A.~De Martino, K.~Diki, and A.~Guzm\'an Ad\'an,
\emph{The Fueter--Sce mapping and the Clifford--Appell polynomials},
2023 preprint, \href{https://arxiv.org/abs/2305.06998}{arXiv:2305.06998}.
The kernel and its relation to generalized Cauchy--Kovalevskaya extension
provide relevant background; no priority is claimed here for the polynomial
kernel.

\bibitem{perotti2024}
A.~Perotti, \emph{A local Cauchy integral formula for slice-regular functions},
Computational Methods and Function Theory \textbf{24} (2024), 185--203.
Published online 30 May 2023.
\href{https://doi.org/10.1007/s40315-023-00485-5}{doi:10.1007/s40315-023-00485-5}.
Theorem~2 states the simply connected component hypothesis used in the
comparison with global inversion.

\end{thebibliography}
\end{document}
